\documentclass[12pt,letterpaper]{article}

\usepackage[utf8]{inputenc}
\usepackage[T1]{fontenc}
\usepackage[spanish,es-tabla]{babel}
\usepackage{geometry}
\usepackage{graphicx}
\usepackage{float}
\usepackage{hyperref}
\usepackage{xcolor}
\usepackage{longtable}
\usepackage{array}
\usepackage{enumitem}
\usepackage{caption}
\usepackage{fancyhdr}

\geometry{margin=2.3cm}
\setlength{\parskip}{0.7em}
\setlength{\parindent}{0pt}
\setlist[itemize]{leftmargin=1.5em}
\setlist[enumerate]{leftmargin=1.7em}

\definecolor{rcegreen}{HTML}{08786F}
\definecolor{rceblue}{HTML}{2B6CB0}
\definecolor{softgray}{HTML}{F4F7F8}

\hypersetup{
  colorlinks=true,
  linkcolor=rcegreen,
  urlcolor=rceblue,
  pdftitle={Manual de uso RCE CQL},
  pdfauthor={RCE CQL}
}

\pagestyle{fancy}
\fancyhf{}
\lhead{Manual de uso RCE CQL}
\rhead{\thepage}

\newcommand{\manualfigure}[3]{
  \begin{figure}[H]
    \centering
    \includegraphics[width=\linewidth]{figures/#1}
    \caption{#2}
    \label{fig:#3}
  \end{figure}
}

\begin{document}

\begin{titlepage}
  \centering
  \vspace*{2cm}
  {\Huge \textbf{Manual de uso RCE CQL}\par}
  \vspace{0.5cm}
  {\Large Prototipo educativo con HL7 FHIR, CQL, ELM y CDS Hooks\par}
  \vspace{1.5cm}
  {\large Guia para docente y alumnos\par}
  \vfill
  \begin{tabular}{ll}
    Sistema & RCE CQL educativo \\
    Datos & Pacientes sinteticos \\
    Estandares & HL7 FHIR R4, CQL/ELM, CDS Hooks \\
    Uso esperado & Clases, demostraciones y pruebas controladas \\
  \end{tabular}
  \vfill
  {\large \today\par}
\end{titlepage}

\tableofcontents
\newpage

\section{Proposito del manual}

Este manual explica como usar el RCE CQL en una clase. La idea principal es que
el alumno pueda escribir reglas CQL, validarlas, probarlas contra pacientes
sinteticos y observar como aparecen o desaparecen recomendaciones clinicas.

El sistema no usa pacientes reales. Los nombres, fechas y eventos clinicos de
las capturas son datos sinteticos generados para docencia. Por eso se pueden
mostrar en clases sin exponer informacion sensible.

\section{Inicio de uso}

Al comenzar una clase, el docente puede entregar una unica URL del sistema. Cada
alumno abre esa URL en su navegador y trabaja en modo \textbf{Alumno}. No es
necesario crear una cuenta ni iniciar sesion manualmente.

Antes de iniciar la demostracion, revise en el menu lateral que los servicios
\textbf{API}, \textbf{HAPI} y \textbf{CQL} aparezcan como operativos. Si alguno
no aparece operativo, conviene resolver ese problema antes de crear reglas o
editar pacientes.

Flujo general de uso del software:

\begin{enumerate}
  \item Abrir la pantalla de pacientes.
  \item Seleccionar un paciente sintetico.
  \item Revisar su ficha clinica y sus datos disponibles.
  \item Crear o abrir una regla CQL.
  \item Validar la regla y revisar el ELM generado.
  \item Probar la regla con un paciente.
  \item Publicar la regla para que pueda generar cards CDS.
  \item Cambiar datos del paciente y reevaluar.
  \item Revisar Actividad CDS para ver el paso a paso de la evaluacion.
\end{enumerate}

La etiqueta de sandbox que aparece en la barra superior identifica el espacio
privado de ese navegador. Los cambios de pacientes, reglas y cards generadas por
un alumno quedan dentro de ese sandbox y no modifican el trabajo de otros
alumnos.

\section{Pantalla de pacientes}

La pantalla de pacientes es el punto de entrada para explorar los datos
sinteticos. En la parte izquierda se encuentra la navegacion principal:
Pacientes, Reglas CQL y Actividad CDS. En la parte inferior se ve el estado de
los servicios API, HAPI y CQL.

\manualfigure{01-listado-pacientes.png}{Listado principal de pacientes sinteticos. La tabla muestra edad, sexo, condiciones activas, ultimo encuentro y estado CDS.}{listado-pacientes}

El filtro de cohortes permite seleccionar grupos por edad. Esto es util para
demostrar reglas que dependen de la edad, como reglas para ninos, adolescentes,
adultos o adultos mayores.

\manualfigure{02-filtro-cohortes.png}{Filtro de cohortes. Permite navegar pacientes por grupo etario antes de probar una regla CQL.}{filtro-cohortes}

\section{Ficha clinica del paciente}

Al abrir un paciente, el RCE muestra una ficha con datos demograficos,
resumen clinico, linea temporal y recomendaciones CDS.

\manualfigure{03-ficha-resumen-card.png}{Ficha de paciente con una card CDS activa. La recomendacion aparece porque una regla CQL publicada aplica al paciente.}{ficha-resumen-card}

La ficha se organiza en pestanas. Cada pestana agrupa recursos FHIR
relacionados. El alumno no necesita conocer JSON FHIR para entender los datos:
la interfaz muestra tablas clinicas simples.

\manualfigure{04-ficha-condiciones.png}{Pestana de condiciones. Estas filas representan recursos FHIR \texttt{Condition}.}{ficha-condiciones}

Las observaciones representan mediciones y resultados. En FHIR se modelan como
\texttt{Observation}. Este recurso es clave para reglas que dependen de
signos vitales o laboratorio.

\manualfigure{05-ficha-observaciones.png}{Pestana de observaciones. Incluye peso, talla, frecuencia cardiaca, resultados y otras mediciones FHIR.}{ficha-observaciones}

Las inmunizaciones se muestran como recursos \texttt{Immunization}. En una
clase se pueden usar para explicar que FHIR no solo representa diagnosticos,
sino tambien eventos preventivos como vacunas.

\manualfigure{06-ficha-inmunizaciones.png}{Pestana de inmunizaciones. Cada fila corresponde a una vacuna registrada como recurso FHIR.}{ficha-inmunizaciones}

\section{Editar datos del paciente}

El boton \textbf{Editar dato} abre un formulario controlado. Este formulario
permite modificar datos utiles para probar reglas, pero sin pedir al alumno que
edite JSON FHIR manualmente.

\manualfigure{07-editar-demografia-observaciones.png}{Formulario de edicion. Permite modificar demografia, signos vitales y laboratorio dentro del sandbox.}{editar-demografia-observaciones}

Las modificaciones se guardan como overlay del sandbox. En otras palabras, el
paciente base de HAPI se mantiene como referencia, y el alumno trabaja con una
copia pedagogica de sus cambios. Esto evita que una prueba afecte a otros
alumnos.

\subsection{Recursos clinicos controlados}

Ademas de datos simples, la maqueta permite agregar recursos clinicos
controlados. El menu evita errores de escritura y mantiene las opciones dentro
de un conjunto seguro para docencia.

\manualfigure{08-agregar-recurso-menu.png}{Menu de recursos clinicos del sandbox. Permite agregar condiciones, observaciones, medicamentos, alergias, encuentros, procedimientos, inmunizaciones y ordenes.}{agregar-recurso-menu}

Por ejemplo, se puede agregar un medicamento activo. En FHIR, un medicamento
indicado se representa como \texttt{MedicationRequest}. Esta accion puede ser
usada por una regla CQL para detectar si el paciente tiene un tratamiento
activo.

\manualfigure{09-agregar-medicamento.png}{Ejemplo de medicamento agregado al sandbox. La interfaz usa selects para reducir errores del alumno.}{agregar-medicamento}

La misma lista de recursos aparece cuando el usuario desea agregar otros tipos
clinicos. Esta captura refuerza que el sandbox trabaja sobre recursos FHIR,
pero los presenta con formularios simples.

\manualfigure{11-menu-recursos-clinicos.png}{Selector de recursos clinicos. El alumno elige el tipo de dato que quiere agregar sin escribir JSON.}{menu-recursos-clinicos}

\section{Catalogo de reglas CQL}

La pantalla de reglas muestra el catalogo de reglas disponibles en el sandbox.
Cada regla tiene nombre, version automatica, estado, hook, activacion y alcance.

\manualfigure{10-catalogo-reglas.png}{Catalogo de reglas CQL. La version es automatica y el alcance indica si la regla pertenece al sandbox.}{catalogo-reglas}

Los estados mas importantes para la clase son:

\begin{itemize}
  \item \textbf{Borrador}: la regla esta en edicion.
  \item \textbf{Validada}: el CQL fue traducido correctamente a ELM.
  \item \textbf{Publicada}: la regla puede participar en la evaluacion CDS.
  \item \textbf{Activa}: la regla publicada se evalua cuando ocurre el hook.
\end{itemize}

El versionado es automatico. El alumno no escribe manualmente la version. El
backend mantiene la coherencia entre CQL, ELM y la \texttt{Library} FHIR.

\section{Editor de reglas CQL}

El editor permite escribir CQL directamente dentro del RCE. La parte izquierda
contiene el codigo, y la parte derecha contiene metadata de la regla.

\manualfigure{12-editor-regla-cql.png}{Editor de regla CQL con version automatica. La expresion booleana indica que resultado se usara para decidir si la card aparece.}{editor-regla-cql}

Un ejemplo simple de regla por edad es:

\begin{verbatim}
library ReglaAdultoDemo version '0.1.0'

using FHIR version '4.0.1'

context Patient

define "Aplica":
  AgeInYears() >= 18
\end{verbatim}

La regla anterior se lee asi: para el paciente actual, la expresion
\texttt{Aplica} devuelve verdadero si la edad en anos es mayor o igual a 18.
Si devuelve verdadero y la regla esta publicada y activa, el RCE muestra una
card CDS.

\subsection{Validar y ver ELM}

El boton \textbf{Validar} envia el CQL al servicio de traduccion. Si el CQL es
correcto, el sistema muestra el ELM generado. Esto permite ensenar que CQL es
el lenguaje que escribe el alumno, mientras que ELM es la representacion que se
puede ejecutar.

\manualfigure{13-elm-validado.png}{ELM generado tras validar la regla. El mensaje indica que la regla esta lista para publicar.}{elm-validado}

\subsection{Publicar regla}

Publicar una regla significa dejarla disponible para que participe en la
evaluacion CDS. En este paso el backend asigna la siguiente version automatica
y guarda la regla como una \texttt{Library} FHIR.

\manualfigure{14-publicar-regla.png}{Dialogo de publicacion. Muestra la Library, ELM guardado, version automatica y alcance del sandbox.}{publicar-regla}

\section{Probar reglas con pacientes}

Antes o despues de publicar, el alumno puede probar la regla contra un paciente
seleccionado. La prueba muestra si la regla aplica o no aplica y que recursos
FHIR fueron usados.

\manualfigure{16-prueba-no-aplica.png}{Prueba de regla contra un paciente menor de edad. La regla no aplica porque la condicion CQL devuelve falso.}{prueba-no-aplica}

Este ejemplo es pedagogicamente importante: no toda regla debe generar una
card. Si el resultado CQL es falso, el comportamiento correcto es no mostrar
recomendaciones.

\section{Cards CDS en la ficha}

Cuando una regla publicada y activa devuelve verdadero, aparece una card CDS en
la ficha del paciente. La card muestra severidad, resumen, detalle y fuente de
la regla.

\manualfigure{15-paciente-cards-publicadas.png}{Paciente adulto con cards CDS activas. Las recomendaciones fueron generadas por reglas CQL publicadas.}{paciente-cards-publicadas}

La card no cambia datos del paciente por si sola. Su funcion es recomendar o
avisar. Si en el futuro una card incluye una sugerencia, el usuario debe
confirmarla antes de aplicar cambios.

\section{Demostracion: cambiar edad para activar o desactivar una regla}

Una demostracion facil para clases es usar una regla de edad. Primero se toma
un paciente adulto que cumple la condicion y luego se cambia su fecha de
nacimiento para que deje de cumplirla.

\manualfigure{17-editar-fecha-nacimiento.png}{Cambio de fecha de nacimiento en el sandbox. Al modificar la fecha cambia la edad efectiva del paciente.}{editar-fecha-nacimiento}

Despues de guardar, el backend reevalua las reglas. Si el paciente queda menor
de 18 anos, la regla de adulto ya no aplica y la card desaparece.

\manualfigure{18-paciente-sin-cards.png}{Paciente despues del cambio de edad. No hay recomendaciones activas porque la condicion CQL devolvio falso.}{paciente-sin-cards}

Este flujo ayuda a explicar tres ideas:

\begin{enumerate}
  \item Los datos clinicos del paciente se representan como HL7 FHIR.
  \item La regla CQL consulta esos datos y devuelve verdadero o falso.
  \item CDS Hooks muestra cards solo cuando hay recomendaciones aplicables.
\end{enumerate}

\section{Actividad CDS}

La pantalla Actividad CDS muestra el historial de evaluaciones. No es una tabla
de pacientes ni una lista de errores: es una vista para explicar que evaluacion
ocurrio, con que paciente, en que momento CDS Hooks y cuantas cards se
generaron.

\manualfigure{19-actividad-cds.png}{Actividad CDS. Cada fila representa una evaluacion de reglas CQL contra datos FHIR del sandbox.}{actividad-cds}

Al presionar \textbf{Paso a paso}, se abre una trazabilidad pensada para la
clase. La interfaz evita poner como foco identificadores tecnicos y privilegia
el flujo conceptual.

\manualfigure{20-paso-a-paso-cds-hooks.png}{Paso a paso CDS Hooks. Muestra el flujo desde el momento clinico hasta el resultado de la card.}{paso-a-paso-cds-hooks}

La lectura recomendada para explicar esta pantalla es:

\begin{enumerate}
  \item El alumno abre o reevalua una ficha.
  \item Ese momento se representa como un hook, por ejemplo \texttt{patient-view}.
  \item El backend toma el paciente y el sandbox del navegador.
  \item Los datos se leen como recursos HL7 FHIR.
  \item Las reglas CQL activas se evaluan sobre esos datos.
  \item Si una regla devuelve verdadero, aparece una card CDS.
\end{enumerate}

\section{Rol Alumno y rol Docente}

El modo Alumno permite explorar pacientes, crear reglas, validar CQL, probar
reglas y publicar reglas dentro de su propio sandbox. Esto permite que cada
alumno trabaje sin afectar a los demas.

El modo Docente esta pensado para preparar reglas compartidas y conducir la
clase. En una instalacion real, el acceso docente debe protegerse con clave o
autenticacion institucional.

\section{Guion corto para una clase}

\begin{enumerate}
  \item Abrir la pantalla de pacientes y explicar que los datos son sinteticos.
  \item Filtrar por cohorte para elegir un paciente adulto.
  \item Abrir la ficha y mostrar que los datos se agrupan como recursos FHIR.
  \item Ir a Reglas CQL y abrir una regla por edad.
  \item Explicar la expresion \texttt{AgeInYears() >= 18}.
  \item Validar la regla y mostrar el ELM.
  \item Publicar la regla.
  \item Volver al paciente y reevaluar para ver la card.
  \item Editar la fecha de nacimiento para dejar al paciente menor de edad.
  \item Guardar y mostrar que la card desaparece.
  \item Abrir Actividad CDS y revisar el paso a paso.
\end{enumerate}

\section{Buenas practicas de uso}

\begin{itemize}
  \item Usar siempre pacientes sinteticos para demostraciones.
  \item Explicar que HAPI FHIR es la fuente de datos y artefactos, no una base
  SQL consultada directamente por el alumno.
  \item Recordar que CQL no escribe datos por si solo: evalua condiciones.
  \item Validar antes de publicar para asegurar que existe ELM.
  \item Usar el sandbox para que cada alumno pruebe cambios sin afectar a otros.
  \item Leer una ausencia de cards como un resultado valido: la regla no aplica.
\end{itemize}

\section{Glosario rapido}

\begin{longtable}{>{\bfseries}p{0.24\linewidth}p{0.68\linewidth}}
Termino & Explicacion simple \\
\hline
RCE & Aplicacion que simula una ficha clinica educativa. \\
FHIR & Estandar HL7 para representar datos clinicos como recursos. \\
HAPI FHIR & Servidor usado para guardar y consultar recursos FHIR. \\
CQL & Lenguaje para escribir reglas clinicas. \\
ELM & Resultado traducido de CQL, listo para evaluacion. \\
CDS Hooks & Forma estandar de pedir recomendaciones clinicas en un punto del flujo. \\
Card CDS & Recomendacion visible para el usuario. \\
Sandbox & Espacio privado de pruebas por navegador. \\
Library & Recurso FHIR que guarda CQL, ELM y metadata de una regla. \\
\end{longtable}

\end{document}
